% GENERATED FILE. Edit canonical lessons, then run npm run generate:books.
% canonical-source-hash: fnv1a64:5338514844fba6d4
% canonical-lessons: GE-C296-geburt, GE-C296-schwanger, GE-C296-schwangerschaft, GE-C296-taufen, GE-C296-taufe

\chapter{Die Geburt}
\label{ch:ge-die-geburt}
\hlchaptermodality{296}

\begin{chapteropening}
\textbf{By the end of this chapter:} \emph{I can tell how a life begins: birth, pregnant, pregnancy, to baptise and a christening.}

Complete the last lesson of chapter 296: I can tell how a life begins: birth, pregnant, pregnancy, to baptise and a christening.
\end{chapteropening}

\section[die Geburt]{die Geburt — the birth}

\label{lesson:GE-C296-geburt}



\begin{quote}
Before the new one: say the German for to damage, then the German for to destroy.
\end{quote}

\subsection*{You'll want to know: die Geburt}
\textbf{die Geburt} — \textquotedblleft{}the birth\textquotedblright{}.

\begin{quote}
Bei der Geburt wog das Baby drei Kilo.
\end{quote}

\begin{quote}
\emph{At birth the baby weighed three kilos.}
\end{quote}

Plural: \emph{die Geburten}. \emph{Geburtstag} — birthday — is built on it.

\begin{cousinweb}
\emph{Geburt} is inherited from Old High German \emph{giburt}, from Proto-Germanic \emph{gaburþiz}, built on \emph{beraną} \textquotedblleft{}to carry, to bear\textquotedblright{} (EtymDB).
\end{cousinweb}

\begin{grammarlens}[title={What you've built}]
One more word for telling, explaining and describing at B1.
\end{grammarlens}

\subsection*{Your turn}
\begin{itemize}
\raggedright
  \item \emph{Say it:} \emph{die Geburt}
  \item \emph{Say it:} \emph{die Geburt}, once more
  \item \emph{Say it:} a sentence of your own with \emph{die Geburt}
  \item \emph{Recall:} say the German for to damage, then the German for to destroy, then say \emph{die Geburt} again
\end{itemize}

\begin{tcolorbox}[breakable,colback=teal!4,colframe=teal!35!black,title={Before you move on}]
What does die Geburt mean? (\textquotedblleft{}The birth\textquotedblright{}.) Say it once more.
\end{tcolorbox}
\section[schwanger]{schwanger — pregnant}

\label{lesson:GE-C296-schwanger}



\begin{quote}
Before the new one: say the German for to destroy, then the German for the birth.
\end{quote}

\subsection*{You'll want to know: schwanger}
\textbf{schwanger} — \textquotedblleft{}pregnant\textquotedblright{}.

\begin{quote}
Meine Schwester war im Sommer schwanger.
\end{quote}

\begin{quote}
\emph{My sister was pregnant in the summer.}
\end{quote}

\emph{Schwanger werden} — to get pregnant.

\begin{cousinweb}
\emph{schwanger} is inherited from Middle High German \emph{swanger}, Old High German \emph{swangar} (EtymDB); its origin beyond that is not settled.
\end{cousinweb}

\begin{grammarlens}[title={What you've built}]
One more word for telling, explaining and describing at B1.
\end{grammarlens}

\subsection*{Your turn}
\begin{itemize}
\raggedright
  \item \emph{Say it:} \emph{schwanger}
  \item \emph{Say it:} \emph{schwanger}, once more
  \item \emph{Say it:} a sentence of your own with \emph{schwanger}
  \item \emph{Recall:} say the German for to destroy, then the German for the birth, then say \emph{schwanger} again
\end{itemize}

\begin{tcolorbox}[breakable,colback=teal!4,colframe=teal!35!black,title={Before you move on}]
What does schwanger mean? (\textquotedblleft{}Pregnant\textquotedblright{}.) Say it once more.
\end{tcolorbox}
\section[die Schwangerschaft]{die Schwangerschaft — the pregnancy}

\label{lesson:GE-C296-schwangerschaft}



\begin{quote}
Before the new one: say the German for the birth, then the German for pregnant.
\end{quote}

\subsection*{You'll want to know: die Schwangerschaft}
\textbf{die Schwangerschaft} — \textquotedblleft{}the pregnancy\textquotedblright{}.

\begin{quote}
Während der Schwangerschaft arbeitete sie weiter.
\end{quote}

\begin{quote}
\emph{During the pregnancy she went on working.}
\end{quote}

Plural: \emph{die Schwangerschaften}.

\begin{cousinweb}
\emph{Schwangerschaft} is built on \emph{schwanger} \textquotedblleft{}pregnant\textquotedblright{}, Old High German \emph{swangar} (EtymDB), with the noun ending \emph{-schaft}.
\end{cousinweb}

\begin{grammarlens}[title={What you've built}]
One more word for telling, explaining and describing at B1.
\end{grammarlens}

\subsection*{Your turn}
\begin{itemize}
\raggedright
  \item \emph{Say it:} \emph{die Schwangerschaft}
  \item \emph{Say it:} \emph{die Schwangerschaft}, once more
  \item \emph{Say it:} a sentence of your own with \emph{die Schwangerschaft}
  \item \emph{Recall:} say the German for the birth, then the German for pregnant, then say \emph{die Schwangerschaft} again
\end{itemize}

\begin{tcolorbox}[breakable,colback=teal!4,colframe=teal!35!black,title={Before you move on}]
What does die Schwangerschaft mean? (\textquotedblleft{}The pregnancy\textquotedblright{}.) Say it once more.
\end{tcolorbox}
\section[taufen]{taufen — to baptise, to christen}

\label{lesson:GE-C296-taufen}



\begin{quote}
Before the new one: say the German for pregnant, then the German for the pregnancy.
\end{quote}

\subsection*{You'll want to know: taufen}
\textbf{taufen} — \textquotedblleft{}to baptise, to christen\textquotedblright{}.

\begin{quote}
Das Kind wurde in der alten Kirche getauft.
\end{quote}

\begin{quote}
\emph{The child was baptised in the old church.}
\end{quote}

Also to name a ship: \emph{ein Schiff taufen}.

\begin{cousinweb}
\emph{taufen} is inherited from Old High German \emph{toufen}, from Proto-Germanic \emph{daupijaną} \textquotedblleft{}to dip\textquotedblright{}, from a root meaning \textquotedblleft{}deep\textquotedblright{} (EtymDB): baptism by dipping.
\end{cousinweb}

\begin{grammarlens}[title={What you've built}]
One more word for telling, explaining and describing at B1.
\end{grammarlens}

\subsection*{Your turn}
\begin{itemize}
\raggedright
  \item \emph{Say it:} \emph{taufen}
  \item \emph{Say it:} \emph{taufen}, once more
  \item \emph{Say it:} a sentence of your own with \emph{taufen}
  \item \emph{Recall:} say the German for pregnant, then the German for the pregnancy, then say \emph{taufen} again
\end{itemize}

\begin{tcolorbox}[breakable,colback=teal!4,colframe=teal!35!black,title={Before you move on}]
What does taufen mean? (\textquotedblleft{}To baptise, to christen\textquotedblright{}.) Say it once more.
\end{tcolorbox}
\section[die Taufe]{die Taufe — the christening, the baptism}

\label{lesson:GE-C296-taufe}



\begin{quote}
Before the new one: say the German for the pregnancy, then the German for to baptise.
\end{quote}

\subsection*{You'll want to know: die Taufe}
\textbf{die Taufe} — \textquotedblleft{}the christening, the baptism\textquotedblright{}.

\begin{quote}
Nach der Taufe gab es Kuchen für alle.
\end{quote}

\begin{quote}
\emph{After the christening there was cake for everyone.}
\end{quote}

Plural: \emph{die Taufen}.

\begin{cousinweb}
\emph{Taufe} is built on \emph{taufen} \textquotedblleft{}to baptise\textquotedblright{}, from Proto-Germanic \emph{daupijaną} \textquotedblleft{}to dip\textquotedblright{} (EtymDB).
\end{cousinweb}

\begin{grammarlens}[title={What you've built}]
One more word for telling, explaining and describing at B1.
\end{grammarlens}

\subsection*{Your turn}
\begin{itemize}
\raggedright
  \item \emph{Say it:} \emph{die Taufe}
  \item \emph{Say it:} \emph{die Taufe}, once more
  \item \emph{Say it:} a sentence of your own with \emph{die Taufe}
  \item \emph{Recall:} say the German for the pregnancy, then the German for to baptise, then say \emph{die Taufe} again
\end{itemize}

\begin{tcolorbox}[breakable,colback=teal!4,colframe=teal!35!black,title={Before you move on}]
What does die Taufe mean? (\textquotedblleft{}The christening, the baptism\textquotedblright{}.) Say it once more.
\end{tcolorbox}
